\documentclass[11pt]{article}
\usepackage{fontspec,amsmath,amssymb,amsthm,hyperref}
\usepackage[margin=1in]{geometry}
\setlength{\emergencystretch}{3em}
\title{Statistical target, initialization and Gaussian innovations}
\author{LangevinLean mathematical exposition}
\date{}
\begin{document}\maketitle
\noindent Source module: \texttt{IPLA/Model.lean}. Compile this proof document with XeLaTeX or LuaLaTeX. The accompanying HTML pairs the mathematical argument with the exact Lean statements and records the scope of the formalization.


This file names the original parameter–latent statistical objects and the law of the actual algorithm. Definitions identify a density and a candidate maximizer; normalization, existence and concentration are proved in later modules.


\section{The marginal likelihood and its maximizer}


The pairing map forms \(\operatorname{pair}(\theta,x)=(\theta,x)\) in the orthogonal joint space. The unnormalized parameter marginal is

\[M(\theta)=\int_{\mathbb R^{d_x}}e^{-U(\theta,x)}\,dx.\]

The definition uses the nonnegative extended integral, so it remains meaningful before finiteness has been established. A marginal maximizer is a point \(\theta_*\) such that \(M(\theta)\le M(\theta_*)\) for every \(\theta\). If \(U=-\log p_\theta(x,y)\), this is maximum marginal likelihood; a prior included in \(U\) gives the corresponding MAP interpretation. The deprecated name is an abbreviation for exactly the same predicate.




\textbf{Correspondence with the displayed declarations.}


\begin{itemize}

\item \texttt{pair}: The original parameter/latent pair \texttt{(theta,x)}, parameter first. For fixed observations \texttt{y}, the standard IPLA potential is \texttt{U y (pair theta x) = -log p\_theta(x,y)}. The geometry and dynamics work with the fixed slice \texttt{U y}.
\item \texttt{marginalDensity}: Unnormalized marginal density of the parameter, integrating out the latent variable. For \texttt{U\_y(theta,x) = -log p\_theta(x,y)}, this is the marginal likelihood \texttt{p\_theta(y)}. A parameter prior may instead be included in \texttt{U\_y}.
\item \texttt{IsMarginalMaximizer}: The parameter maximizes the integrated joint density. This is maximum likelihood in standard IPLA notation, and MAP when the parameter prior is included in \texttt{U}. Its existence is not asserted by this definition.
\item \texttt{IsMarginalMAP}: Compatibility name for the prior-inclusive formulation.
\end{itemize}


\noindent Lean declarations: \texttt{pair}, \texttt{marginalDensity}, \texttt{IsMarginalMaximizer}, \texttt{IsMarginalMAP}.

\section{The actual Gibbs target and all Gaussian assumptions}


For \(V_N\) defined by the geometry module, the target is \(\Pi_N=Z_N^{-1}e^{-V_N}\,dz\), with \(Z_N=\int e^{-V_N}\,dz\). Both normalization and density are initially extended-valued; this definition alone asserts no probability property. The Gaussian innovation structure asks that every coordinate \(\xi_{k,a}\) have law \(\mathcal N(0,1)\), that all coordinates at all times be mutually independent, and that the entire noise path be independent of the initialization. These are explicit distributional hypotheses on the supplied random variables.




The Gaussian structure concerns the entire countable collection \((k,a)\mapsto\xi_{k,a}\), not only coordinate independence at one time. The separate initialization field states independence of \(Z_0\) from the whole innovation path. Together these conditions imply that every transition receives a fresh full state-space standard Gaussian. They are the probabilistic assumptions later used to equate a supplied implementation with the canonical Markov laws; moment membership in \(L^2\) alone would not establish either the Gaussian distribution or independence.




\textbf{Correspondence with the displayed declarations.}


\begin{itemize}

\item \texttt{gibbsTarget}: The normalized density in (49), with \texttt{U\_N} the actual lifted potential. The integral and scalar are extended-valued. \texttt{gibbsTarget\_isProbability} derives positive finite normalization from the original joint strong convexity, before any conversion to real numbers.
\item \texttt{GaussianInnovations}: The actual update noise is independent standard Gaussian in every coordinate and time, and independent of initialization. These are distributional hypotheses, not a claim to have constructed a random sequence.
\end{itemize}


\noindent Lean declarations: \texttt{gibbsTarget}, \texttt{GaussianInnovations}.

\section{The law of the L\ensuremath{{}^{2}} recursion}


A square-integrable initial state and square-integrable noise vectors lie in \(L^2(P)\). With \(h=\gamma/N\), define the random sequence through the shared ULA recursion

\[Z_{k+1}=Z_k-h\nabla V_N(Z_k)+\sqrt{2h}\,\xi_k,\qquad \rho_k=(Z_k)_\#P.\]

The Lipschitz estimate for \(\nabla V_N\) ensures the drift update stays in \(L^2\). The definition of \texttt{laws} is precisely this pushforward. The dynamics module proves that its representatives agree almost everywhere with the ordinary pointwise IPLA recursion, so the eventual theorem concerns this actual algorithm.




\textbf{Correspondence with the displayed declarations.}


\begin{itemize}

\item \texttt{laws}: The laws of the actual noise-driven IPLA recursion, in scaled coordinates. The coordinate identities in \texttt{Geometry.lean} show that this is precisely the usual average-gradient parameter step and ordinary particle Langevin steps.
\end{itemize}


\noindent Lean declarations: \texttt{laws}.

\end{document}
